\begin{definition}[Bounded tagged entries]
\label{def:p04-tagged-history-entry}
The following are eliminable abbreviations for arithmetic formulas:
\[
\begin{aligned}
 \mathsf E(s,t,j,v)&\defeq
   \exists i<S(t+j)\times S(t+j)\,
   \bigl(\mathsf P(t,j,i)\land\mathsf B(s,i,v)\bigr),\\
 \mathsf E^{(2)}(s,t,a,b,v)&\defeq
   \exists j<S(a+b)\times S(a+b)\,
   \bigl(\mathsf P(a,b,j)\land\mathsf E(s,t,j,v)\bigr).
\end{aligned}
\]
In each expansion the displayed bound variables, and those internal to
\(\mathsf B\), are renamed to fresh variables before substitution of
the argument terms.  Thus a bound never contains the variable it binds.
All quantifiers are marked bounded; neither pairing nor an array-access
function is added to \(\Lang_A\).  In numerical descriptions below,
write \(s[t,j]=\beta(s,\operatorname{pair}_{\N}(t,j))\) and
\(s[t,a,b]=\beta(s,\operatorname{pair}_{\N}
 (t,\operatorname{pair}_{\N}(a,b)))\).
These bracket expressions denote natural numbers in the metatheory;
their use inside an arithmetic formula is replaced by a fresh value
variable and the corresponding \(\mathsf E\) or \(\mathsf E^{(2)}\)
conjunct.
\end{definition}

\begin{lemma}[Tagged extraction and finite packing]
\label{lem:p04-tagged-history-extraction}
Both formulas in Definition~\ref{def:p04-tagged-history-entry} are
bounded.  In \(\mathcal N\), for each choice of their other numerical
arguments they hold of exactly one value \(v\), respectively
\(s[t,j]\) and \(s[t,a,b]\), and this value is at most \(s\).
Fix for each family tag whether its positions are one- or two-dimensional.
Any finite assignment of numerical values to distinct positions of these
tagged families is realized by some \(s\).  The paired indices themselves are not
required to be at most \(s\).
\textit{Status: complete proof.}
\end{lemma}

\begin{proof}
For \(a,b\in\N\), put \(u=a+b\).  The numerical pairing satisfies
\[
 \operatorname{pair}_{\N}(a,b)
 =\frac{u(u+1)}2+b
 \leq u(u+1)+u
 =u^2+2u<(u+1)^2.
\]
The inequality also holds when \(u=0\).  Consequently the bound in
\(\mathsf E\) includes the unique value
\(i=\operatorname{pair}_{\N}(t,j)\); the graph formula
\(\mathsf P\) excludes all other values.  Lemma~\ref{lem:goedel-beta}(a)
then gives exactly the value \(\beta(s,i)\), bounded by \(s\).
Applying the same pairing bound first to \((a,b)\), and then applying
the just-proved assertion about \(\mathsf E\), proves the claim for
\(\mathsf E^{(2)}\).  Every displayed bound is an arithmetic term,
and substituting the already bounded formulas \(\mathsf P,\mathsf B\)
after fresh renaming introduces only bounded quantifiers.

For finite packing, flatten each two-dimensional position by the inner
pairing, then pair its position with its family tag.  Each tag has one
fixed positional arity.  Within a one-dimensional family distinct positions
are distinct integers; within a two-dimensional family injectivity of the
inner pairing preserves distinctness.  The outer pairing separates the
families, so distinct tagged positions give distinct indices.  If there are no
assigned positions, any \(s\) realizes the empty assignment.  Otherwise
let \(I\) be one more than the greatest resulting index.  Form the
finite sequence of length \(I\) having the assigned value at each
specified index and zero at every other index.  Lemma~\ref{lem:goedel-beta}(b)
gives an \(s\) realizing this sequence.  Its part (a) automatically
bounds every assigned value by \(s\).  The indices were computed from
the finite assignment before choosing \(s\), so this construction
does not make the indices depend on the code subsequently obtained.
\end{proof}

\begin{definition}[History coding]
\label{def:history-coding}
Use the tagged-entry convention of
Definition~\ref{def:p04-tagged-history-entry}, with the following fixed
family tags.  A two-dimensional family uses the paired position \((a,b)\)
in the order displayed; all other families use one position.
\[
\begin{array}{c|l|l}
 \text{tag}&\text{stored values}&\text{positions}\\ \hline
 0&\text{header }(k,c_0,n,b,A,L,T,H,W,\ell,y)&0,\ldots,10\\
 1&\text{description entries }r_i&i<k\\
 2&\text{description decoding tails }c_i&i<k\\
 3&\text{input-alphabet bitmap digits }\eta_j&j<b\\
 4&\text{bitmap tail values }B_j&j\leq b\\
 5&\text{canonical input digits }d_j&j<L\\
 6&\text{input prefix values }v_j&j<L\\
 7&\text{configuration states }q_t&t\leq T\\
 8&\text{shifted head positions }h_t&t\leq T\\
 9&\text{tape symbols }a_{t,j}&(t,j),\ t\leq T,\ j<W\\
 10&\text{tape prefix values }u_{t,j}&(t,j),\ t\leq T,\ j<W\\
 11&\text{complete fixed-width tape codes }z_t&t\leq T\\
 12&\text{canonical output digits }o_j&j<\ell\\
 13&\text{output prefix values }w_j&j<\ell
\end{array}
\]
Here \(n,b,A\) have the meanings fixed by the serialized program;
\(L\geq1\) is the input length, \(T\) the number of transitions,
\(H=T\) the tape origin, and \(W=L+2T+1\) its fixed width.
Column \(j<W\) represents physical cell \(j-H\).  Bitmap tails
satisfy \(B_b=0\), \(B_j=2B_{j+1}+\eta_j\) for \(j<b\),
and \(B_0=A\); this fixes exactly the bitmap of the stated width.
Each tape expansion runs left to right:
\(u_{t,0}=a_{t,0}\),
\(u_{t,j+1}=b\,u_{t,j}+a_{t,j+1}\) for \(j+1<W\), and
\(z_t=u_{t,W-1}\), with \(a_{t,j}<b\).
The output length \(\ell\geq1\) and value \(y\) refer to the
canonical prefix at physical cells \(0,\ldots,\ell-1\), ending
at the first blank cell \(\ell\), so \(H+\ell<W\).
For that prefix \(w_0=o_0\), \(w_{j+1}=2w_j+o_{j+1}\), and
\(w_{\ell-1}=y\).  Both input and output words use the earlier
canonical binary convention, including the single digit for zero and
the absence of leading zeros for positive values.

These are numerical descriptions of the stored data.  Every lookup has
the bounded relational expansion just proved; an index such as
\(\operatorname{pair}_{\N}(t,j)\) is never an extra arithmetic
function symbol.  A number \(s\) is a
\emph{history code} for the pair \((e,x)\) iff its entries comprise, under
this scheme:
\begin{enumerate}[label=(h\arabic*)]
\item the binary digits \(d_0,\ldots,d_{L-1}\) of \(x\), most significant
      digit first, with certified expansion values
      \(v_0,\ldots,v_{L-1}\) satisfying \(v_0=d_0\),
      \(v_{j+1}=2v_j+d_{j+1}\) for \(j+1<L\), and \(v_{L-1}=x\), together
      with the certificates \(v_j\leq x\) and \(d_j\leq1\) for all
      \(j<L\);
\item the exact serialized description of \(e\) from
      Definition~\ref{def:p04-program-serialization}: a decoding chain consisting of
      \(\mathsf P(k,c_0,e)\), \(\mathsf P(r_i,c_{i+1},c_i)\) for
      \(i<k-1\), and \(r_{k-1}=c_{k-1}\), unpacking the description tuple of
      length \(k\) into its entries \(r_0,\ldots,r_{k-1}\), each entry
      interpreted as the state count, tape-alphabet size, input-alphabet
      bitmap, and then the lexicographically ordered five-field row codes.
      Each row is unpacked by four \(\mathsf P\)-steps; the bitmap is
      accompanied by binary expansion certificates.  These data include
      all length, range, alphabet and full-table validity checks of the fixed numbering under its
      designated-slot normalization;
\item configurations \(C_0,\ldots,C_T\) of the decoded machine, each tape
      coded in the tape-alphabet base with certified expansions as in
      (h1), the configuration \(C_0\) initial for the input word certified
      by (h1), and each successor configuration matching its predecessor
      and the decoded table by bounded digit comparisons; and
\item a final configuration whose state is the accepting state and whose
      tape prefix before the first blank is a valid binary numeral in the
      sense of the fixed output convention.
\end{enumerate}
Every stored datum in these clauses is an entry of \(s\), unique and
at most \(s\) by \cref{lem:goedel-beta}(a).  Row projections are
bounded by their stored row code using \cref{lem:pairing-graph}.
The tagged-entry indices have the separate bounds of
\cref{lem:p04-tagged-history-extraction}; the formulas below give
explicit bounds for derived lengths and last-position witnesses.
\end{definition}

\begin{definition}[Bounded input and table formulas]
\label{def:p04-input-table-formulas}
The formulas below give the complete arithmetic checks for clauses
(h1) and (h2) of Definition~\ref{def:history-coding}.  Each displayed
name is an eliminable formula abbreviation, with bound variables renamed
fresh before its arguments are substituted.  Fixed numerical constants
other than \(0\) are written as their numeral terms.  In particular,
neither a tuple constructor nor an entry-access function is added to
\(\Lang_A\).  Use the following bounded comparison and digit formulas:
\[
\begin{aligned}
 \mathsf{Lt}(u,v)&\defeq\exists z<v\,(u+Sz=v),\\
 \mathsf{Le}(u,v)&\defeq\exists z\leq v\,(z+u=v),\\
 \mathsf{Bit}(d)&\defeq(d=0\lor d=\bar1).
\end{aligned}
\]
In \(\mathcal N\), the first two formulas express respectively
\(u<v\) and \(u\leq v\), including when \(v=0\).  These formulas
are used for comparisons inside matrices; all quantifier bounds below
have the marked syntax of Definition~\ref{def:language-arithmetic}
and its following conventions.

For input cells and successive input prefixes, put
\[
\begin{aligned}
 \mathsf I(x,s,L)\defeq{}&
  \forall j<L\,\exists d\leq s\,\exists v\leq s\,
   \bigl(\mathsf E(s,\bar5,j,d)\land
         \mathsf E(s,\bar6,j,v)\land
         \mathsf{Bit}(d)\land\mathsf{Le}(v,x)\bigr),\\
 \mathsf J(s,\lambda)\defeq{}&
  \forall j<\lambda\,\exists v\leq s\,\exists w\leq s\,
  \exists d\leq s\,
   \bigl(\mathsf E(s,\bar6,j,v)\land
         \mathsf E(s,\bar6,Sj,w)\land
         \mathsf E(s,\bar5,Sj,d)\land
         w=\bar2\times v+d\bigr).
\end{aligned}
\]
The variable \(\lambda\) denotes the last input position in this
formula; it is unrelated to the output length \(\ell\) stored in
the history header.  Define
\[
\begin{aligned}
 \mathsf H_1(x,s)\defeq{}&
  \exists L\leq s\,\exists\lambda\leq s\,\exists d\leq s\,
  \bigl[\mathsf E(s,0,\bar5,L)\land L=S\lambda\\
 &\quad\land\mathsf E(s,\bar5,0,d)
       \land\mathsf E(s,\bar6,0,d)
       \land\mathsf E(s,\bar6,\lambda,x)\\
 &\quad\land\bigl((x=0\land L=\bar1\land d=0)
              \lor(\neg(x=0)\land d=\bar1)\bigr)\\
 &\quad\land\mathsf I(x,s,L)\land\mathsf J(s,\lambda)\bigr].
\end{aligned}
\]

The following graph unpacks a five-field row using exactly the
right-nested tuple convention of
Definition~\ref{def:p04-program-serialization}:
\[
\begin{aligned}
 \mathsf P_5(q,a,q',a',d,\rho)\defeq{}&
  \exists v_1\leq\rho\,\exists v_2\leq\rho\,
  \exists v_3\leq\rho\,\bigl[
       \mathsf P(q,v_1,\rho)\land
       \mathsf P(a,v_2,v_1)\\
 &\hspace{35mm}\land\mathsf P(q',v_3,v_2)
       \land\mathsf P(a',d,v_3)\bigr],\\
 \mathsf R_5(s,i,q,a,q',a',d)\defeq{}&
  \exists\rho\leq s\,
   \bigl[\mathsf E(s,\bar1,\bar3+i,\rho)
           \land\mathsf P_5(q,a,q',a',d,\rho)\bigr].
\end{aligned}
\]
Thus \(i\) is the zero-based row number, while \(\bar3+i\) is
its position in the full description.  The row fields need no new
stored family.  The three internal pairing tails are bounded by
\(\rho\); the four pairing equations also force every row field
to be at most \(\rho\).  By itself \(\mathsf R_5\) imposes
neither table membership nor the machine-field ranges: the full check
below applies it only to row numbers \(i<m\) and supplies those ranges.

For a description whose last position is \(\lambda\), let
\[
\begin{aligned}
 \mathsf D(s,\lambda)\defeq{}&
  \bigl[\forall i<\lambda\,\exists r\leq s\,
    \exists c\leq s\,\exists c'\leq s\,
    \bigl(\mathsf E(s,\bar1,i,r)\land
          \mathsf E(s,\bar2,i,c)\\
 &\hspace{32mm}\land\mathsf E(s,\bar2,Si,c')
          \land\mathsf P(r,c',c)\bigr)\bigr]\\
 &\quad\land\exists v\leq s\,
    \bigl(\mathsf E(s,\bar1,\lambda,v)
          \land\mathsf E(s,\bar2,\lambda,v)\bigr).
\end{aligned}
\]
The last conjunct implements the identity coding of a one-entry tail;
there is no extra zero terminator.  The bitmap check is
\[
\begin{aligned}
 \mathsf M(s,b,A)\defeq{}&
  \mathsf E(s,\bar4,0,A)\land\mathsf E(s,\bar4,b,0)\\
 &\land\mathsf E(s,\bar3,0,0)
       \land\mathsf E(s,\bar3,\bar1,\bar1)
       \land\mathsf E(s,\bar3,\bar2,\bar1)\\
 &\land\forall j<b\,\exists d\leq s\,\exists v\leq s\,
      \exists w\leq s\,
      \bigl[\mathsf E(s,\bar3,j,d)\land
            \mathsf E(s,\bar4,j,v)\\
 &\hspace{30mm}\land\mathsf E(s,\bar4,Sj,w)
            \land\mathsf{Bit}(d)\land v=\bar2\times w+d\bigr].
\end{aligned}
\]
For \(\nu\) nonhalting states, the complete ordered table check is
\[
\begin{aligned}
 \mathsf K(s,n,b,\nu)\defeq{}&
  \forall u<\nu\,\forall a<b\,
  \exists q\leq s\,\exists q'\leq s\,
  \exists a'\leq s\,\exists d\leq s\,\bigl[\\
 &\quad\bigl((u=0\land q=0)
           \lor(\neg(u=0)\land q=u+\bar2)\bigr)\\
 &\quad\land\mathsf R_5(s,u\times b+a,q,a,q',a',d)\\
 &\quad\land\mathsf{Lt}(q,n)
           \land\neg(q=\bar1)\land\neg(q=\bar2)\\
 &\quad\land\mathsf{Lt}(q',n)
           \land\mathsf{Lt}(a',b)\land\mathsf{Bit}(d)\bigr].
\end{aligned}
\]
Here \(d=0\) denotes a left move and \(d=\bar1\) a right
move.  The ordered blocks have states \(0,3,4,\ldots,n-1\), with
the tape symbols in increasing order within each block.  Finally set
\[
\begin{aligned}
 \mathsf H_2(e,s)\defeq{}&
  \exists k\leq s\,\exists c\leq s\,\exists n\leq s\,
  \exists b\leq s\,\exists A\leq s\,\exists\nu\leq s\,
  \exists m\leq s\,\exists\lambda\leq s\,\bigl[\\
 &\quad\mathsf E(s,0,0,k)\land\mathsf E(s,0,\bar1,c)
        \land\mathsf E(s,0,\bar2,n)\\
 &\quad\land\mathsf E(s,0,\bar3,b)
        \land\mathsf E(s,0,\bar4,A)\\
 &\quad\land n=\nu+\bar2\land\neg(\nu=0)
        \land\mathsf{Le}(\bar3,b)\\
 &\quad\land m=\nu\times b\land k=\bar3+m
        \land k=S\lambda\land\mathsf P(k,c,e)\\
 &\quad\land\mathsf E(s,\bar2,0,c)
        \land\mathsf E(s,\bar1,0,n)
        \land\mathsf E(s,\bar1,\bar1,b)
        \land\mathsf E(s,\bar1,\bar2,A)\\
 &\quad\land\mathsf D(s,\lambda)
        \land\mathsf M(s,b,A)\land\mathsf K(s,n,b,\nu)\bigr].
\end{aligned}
\]
The two occurrences of \(\lambda\) in \(\mathsf H_1\) and
\(\mathsf H_2\) are separately bound and are renamed apart on
expansion.  The checks concern the exact raw program encoded by \(e\).
They neither change its transition table to an output-cleaning wrapper
nor assert that it accepts the input.  These explicit formulas and their
proofs below are the handbook's construction for its fixed numbering.
\end{definition}

\begin{lemma}[Correctness of the input formula]
\label{lem:p04-input-formula}
The formula \(\mathsf H_1\) is bounded.  For every \(x,s\in\N\),
\(\mathcal N\entails\mathsf H_1(\bar x,\bar s)\) if and only if
the header input length and tagged input digits and prefixes satisfy
exactly clause (h1), with the canonical binary convention of
Definition~\ref{def:history-coding}.  No condition on the other tagged
families is imposed by this equivalence.
\textit{Status: complete proof.}
\end{lemma}

\begin{proof}
Expanding \(\mathsf I,\mathsf J,\mathsf{Bit},\mathsf{Le}\), and
then \(\mathsf E\), introduces only marked bounded quantifiers and
atomic formulas of \(\Lang_A\), combined by Boolean connectives.
Fresh renaming ensures that no bound contains the variable it binds.
This proves boundedness in the stated syntactic convention.

Suppose first that \(\mathsf H_1(x,s)\) holds in \(\mathcal N\).
The unique-extraction lemma identifies the witnesses \(L\) and \(d\)
with the header length and the first stored input digit.  The equation
\(L=S\lambda\) makes the length positive.  Write \(d_j=s[5,j]\)
and \(v_j=s[6,j]\) in the metatheory.  Uniqueness of extraction in
the universal clauses gives, for every indicated position,
\[
 d_j\in\{0,1\},\qquad v_j\leq x,\qquad
 v_0=d_0,\qquad v_{j+1}=2v_j+d_{j+1},\qquad v_{L-1}=x.
\]
If \(x=0\), the explicit first disjunct forces \(L=1\) and
\(d_0=0\); the other disjunct is false.  If \(x>0\), the second
disjunct forces \(d_0=1\).  The prefix recurrence says precisely
that the word has numerical value \(x\), and the first digit excludes
leading zeros.  For completeness, this characterizes the usual binary
word uniquely: if a word has at least two digits, its last recurrence
is the Euclidean decomposition of its value into twice the preceding
prefix plus a remainder in \(\{0,1\}\).  The preceding prefix is
positive and strictly smaller than the value.  A one-digit positive
word has value \(1\), whereas any longer word starting with \(1\)
has value at least \(2\).  Induction on the positive value, using
uniqueness of quotient and remainder on division by \(2\), therefore
recovers all the digits and the length.  Together with the zero case,
these are exactly the canonical words fixed earlier.  This proves (h1).

Conversely suppose the header length and families 5 and 6 have the
data required by (h1).  Set \(\lambda=L-1\) in the metatheory and
take \(d=d_0\).  Every stored value is at most \(s\) by
Lemma~\ref{lem:p04-tagged-history-extraction}; also
\(\lambda<L\leq s\).  Thus all displayed existential bounds
admit these witnesses.  The canonical zero word satisfies the first
disjunct, and every positive canonical word satisfies the second.
The prefix conditions give every conjunct of \(\mathsf J\) and the
last-prefix equation.  The digit and prefix bounds give \(\mathsf I\).
The bounded comparison witnesses exist because their numerical
inequalities hold.  This proves \(\mathsf H_1(x,s)\).
\end{proof}

\begin{lemma}[Correctness of the serialized-table formula]
\label{lem:p04-table-formula}
The formula \(\mathsf H_2\) is bounded.  For every \(e,s\in\N\),
\(\mathcal N\entails\mathsf H_2(\bar e,\bar s)\) if and only if
the header positions \(0,\ldots,4\) and tagged families
\(1,2,3,4\) certify exactly the valid serialized program \(e\)
specified in clause (h2).  In particular, an invalid program index
satisfies \(\mathsf H_2(e,s)\) for no \(s\).  For a certified
table and \(i<m=(n-2)b\), the row relation \(\mathsf R_5\)
extracts the five fields of row \(i\); it is available for subsequent
transition checks with this row bound.
\textit{Status: complete proof.}
\end{lemma}

\begin{proof}
Each helper in Definition~\ref{def:p04-input-table-formulas} has
only marked bounded quantifiers, arithmetic-term bounds, and Boolean
combinations of previously bounded formulas.  In particular,
\(u\times b+a\) and \(\bar3+i\) are legitimate argument terms
for \(\mathsf E\).  Its internal paired indices have their own
polynomial bounds from Lemma~\ref{lem:p04-tagged-history-extraction};
they need not be at most \(s\).  After fresh renaming and expansion,
\(\mathsf H_2\) is therefore bounded in \(\Lang_A\).

First consider \(\mathsf P_5\) in the standard structure.
The pairing-graph lemma makes its four equations equivalent to
\[
 \rho=\operatorname{pair}_{\N}
  \bigl(q,\operatorname{pair}_{\N}
   (a,\operatorname{pair}_{\N}
    (q',\operatorname{pair}_{\N}(a',d)))\bigr).
\]
This is exactly \(\operatorname{code}_5(q,a,q',a',d)\).
Successive pairing projections are unique, so all five fields are
unique.  Conversely, the three internal tails of any such tuple are
at most \(\rho\), by applying the component bounds successively;
hence their witnesses meet the displayed bounds.  The same bounds
put each field at most \(\rho\).  Together with unique tagged
extraction, this proves the stated meaning of \(\mathsf R_5\)
without an extra row-certificate array.

Suppose \(\mathsf H_2(e,s)\) holds.  Its header witnesses are the
unique stored header values.  The equations
\(n=\nu+2\), \(\nu\ne0\), and \(3\leq b\) give
\(n,b\geq3\), and then
\[
 m=(n-2)b\geq3,\qquad k=3+m\geq6,\qquad \lambda=k-1.
\]
Write \(r_i=s[1,i]\) and \(c_i=s[2,i]\) in the metatheory.
The initial-tail conjunct gives \(c_0=c\).  Unique extraction in
\(\mathsf D\) gives
\[
 c_i=\operatorname{pair}_{\N}(r_i,c_{i+1})\quad(i<k-1),
 \qquad c_{k-1}=r_{k-1}.
\]
Backward induction on \(i\) identifies \(c_i\) with the
right-nested code of \((r_i,\ldots,r_{k-1})\), starting with the
identity code of the last entry.  Thus
\(c=\operatorname{code}_k(r_0,\ldots,r_{k-1})\), and the outer
pairing equation is exactly the required code \(e\).  The other
entry conjuncts give \((r_0,r_1,r_2)=(n,b,A)\).  This also excludes
short descriptions, wrong outer lengths, and extra counted rows:
the specified tuple has exactly \(3+(n-2)b\) entries.  Tagged values
at positions outside the prescribed ranges are irrelevant.

Put \(\eta_j=s[3,j]\) and \(B_j=s[4,j]\).  The bitmap formula
gives \(B_b=0\), \(B_0=A\), and
\(B_j=2B_{j+1}+\eta_j\), with each \(\eta_j\) a bit.
Backward induction gives the numerical identity
\[
 B_j=\sum_{a=j}^{b-1}\eta_a2^{a-j}\quad(0\leq j\leq b),
\]
where the sum is empty for \(j=b\).  This displayed sum and its
powers are metatheoretic statements in the correctness proof, not
terms of the arithmetic formula.  Hence \(A\) has precisely the
width-\(b\) bitmap required by the serialization, with no higher
bits.  The three specified digits exclude the blank symbol and include
both binary input symbols.  They are positions within this bitmap
because \(b\geq3\).

It remains to verify the complete table, rather than just the ranges
of individual rows.  Since \(b>0\), division by \(b\) writes every
integer \(i<m=\nu b\) uniquely as \(i=ub+a\), with
\(u<\nu\) and \(a<b\).  For \(u=0\), the formula \(\mathsf K\)
forces the source state \(q=0\); for \(0<u<\nu=n-2\), it forces
\(q=u+2\).  These states are exactly \(0,3,\ldots,n-1\), in
that order, including the case \(n=3\) when there is just the
state-\(0\) block.  Within each block \(a\) runs once through
\(0,\ldots,b-1\).  The row relation at position \(3+ub+a\)
therefore forces the prescribed lexicographic key, and its remaining
conjuncts force a next state below \(n\), a written symbol below
\(b\), and one of the two directions.  The source state is neither
accepting nor rejecting.  Since pairing projections are unique, each
key has exactly one decoded transition.  Every nonhalting key occurs
and there are no additional rows in the tuple.  The described table
is consequently a total transition function on precisely the required
nonhalting keys.  This proves all of clause (h2).

Conversely, suppose the indicated tags certify the valid serialized
program \(e\).  Take its actual \(k,c_0,n,b,A\) as header
witnesses, and set \(\nu=n-2\), \(m=(n-2)b=k-3\), and
\(\lambda=k-1\) in the metatheory.  The first five witnesses are
stored values, hence at most \(s\).  Also
\(\nu\leq n\leq s\), \(m\leq k\leq s\), and
\(\lambda\leq k\leq s\); no subtraction function is needed in
the formula, whose equations express these relationships.
The actual decoding tails and entries satisfy \(\mathsf D\)
and its bounds by tagged extraction.  The bitmap digits and tail
values satisfy \(\mathsf M\), including its terminal zero and
three distinguished digits.  For each pair \((u,a)\) in the table
check, use the five fields of its actual row.  Its code is an entry
\(\rho\leq s\), and all fields and internal tails are at most
\(\rho\), as proved above.  Thus every existential witness in
\(\mathsf K\) and \(\mathsf P_5\) is admitted by its bound.
The source key, field ranges, and direction satisfy the remaining
conjuncts by validity of the serialization.  All header and outer
pairing conjuncts hold as well, proving \(\mathsf H_2(e,s)\).
The forward implication already shows that no invalid index can
have such a certificate.
\end{proof}

\begin{lemma}[Simultaneous input and table certificates]
\label{lem:p04-input-table-packing}
For every valid program index \(e\) and every \(x\in\N\),
there is an \(s\) such that
\(\mathcal N\entails\mathsf H_1(\bar x,\bar s)
 \land\mathsf H_2(\bar e,\bar s)\).
More generally, any finite assignment agreeing with these certificates
at every shared tagged position can be packed together with them.
\textit{Status: complete proof.}
\end{lemma}

\begin{proof}
Write the canonical word for \(x\) and form its successive numerical
prefixes.  They satisfy the binary recurrence.  Their values are
nondecreasing, since \(2v+d\geq v\) for \(v\geq0\) and a bit
\(d\); thus every prefix is at most the final value \(x\).
Take the actual finite serialized description of \(e\), its right-nested
decoding tails, and its bitmap digits.  Define the bitmap tails backward
from \(B_b=0\) by \(B_j=2B_{j+1}+\eta_j\); the serialization
identity gives \(B_0=A\).  Assign these finitely many values to the
specified header positions and families \(1,\ldots,6\).
The repeated descriptions of a header value or first or last prefix
agree by construction.  The input and table families are otherwise
disjoint.  Add any further finite assignments that agree at overlapping
positions.  Lemma~\ref{lem:p04-tagged-history-extraction} supplies
one \(s\) realizing the assignment, and automatically bounds all its
assigned values by \(s\).  The two preceding correctness lemmas give
the asserted formulas.  This construction does not require an accepting
run: input and table certificates also exist when the program rejects
or diverges.  Clauses (h3) and (h4) are still needed for a computation
certificate.
\end{proof}

\begin{definition}[Bounded configuration, transition and output checks]
\label{def:p04-computation-formula}
Continue the fresh-variable and eliminable-abbreviation conventions of
Definition~\ref{def:p04-input-table-formulas}.  For a time \(t\), tape
width \(W\), and last column \(\mu\), define
\[
\begin{aligned}
 \mathsf{Cfg}(s,t,n,b,W,\mu)\defeq{}&
 \bigl[\exists q\leq s\,\exists h\leq s\,
  (\mathsf E(s,\bar7,t,q)\land\mathsf E(s,\bar8,t,h)
       \land\mathsf{Lt}(q,n)\land\mathsf{Lt}(h,W))\bigr]\\
 &\land\bigl[\forall j<W\,\exists a\leq s\,
    (\mathsf E^{(2)}(s,\bar9,t,j,a)\land\mathsf{Lt}(a,b))\bigr]\\
 &\land\bigl[\exists a\leq s\,
    (\mathsf E^{(2)}(s,\bar9,t,0,a)
       \land\mathsf E^{(2)}(s,\overline{10},t,0,a))\bigr]\\
 &\land\bigl[\forall j<\mu\,\exists u\leq s\,\exists v\leq s\,
       \exists a\leq s\,(
       \mathsf E^{(2)}(s,\overline{10},t,j,u)\\
 &\hspace{24mm}\land\mathsf E^{(2)}(s,\overline{10},t,Sj,v)
       \land\mathsf E^{(2)}(s,\bar9,t,Sj,a)
       \land v=b\times u+a)\bigr]\\
 &\land\exists z\leq s\,
   (\mathsf E(s,\overline{11},t,z)
       \land\mathsf E^{(2)}(s,\overline{10},t,\mu,z)).
\end{aligned}
\]
The initial-tape check converts numerical bits \(0,1\) to their
normalized tape-symbol numbers \(1,2\):
\[
\begin{aligned}
 \mathsf{Init}(s,L,H,W)\defeq{}&
   \mathsf E(s,\bar7,0,0)\land\mathsf E(s,\bar8,0,H)\\
 &\land\bigl[\forall i<L\,\exists d\leq s\,
   (\mathsf E(s,\bar5,i,d)
       \land\mathsf E^{(2)}(s,\bar9,0,H+i,Sd))\bigr]\\
 &\land\forall j<W\,
   \bigl[(\mathsf{Lt}(j,H)\lor\mathsf{Le}(H+L,j))
           \to\mathsf E^{(2)}(s,\bar9,0,j,0)\bigr].
\end{aligned}
\]
For one transition use the following bounded formula:
\[
\begin{aligned}
 \mathsf{Step}(s,t,m,W)\defeq{}&
  \exists q\leq s\,\exists q'\leq s\,
  \exists h\leq s\,\exists h'\leq s\,
  \exists a\leq s\,\exists a'\leq s\,
  \exists d\leq s\,\exists i<m\,\bigl[\\
 &\quad\mathsf E(s,\bar7,t,q)\land\mathsf E(s,\bar7,St,q')
       \land\mathsf E(s,\bar8,t,h)\land\mathsf E(s,\bar8,St,h')\\
 &\quad\land\mathsf E^{(2)}(s,\bar9,t,h,a)
       \land\mathsf R_5(s,i,q,a,q',a',d)
       \land\neg(q=\bar1)\land\neg(q=\bar2)\\
 &\quad\land\bigl((d=0\land h=Sh')
                    \lor(d=\bar1\land h'=Sh)\bigr)
       \land\mathsf E^{(2)}(s,\bar9,St,h,a')\\
 &\quad\land\forall j<W\,
    \bigl[\neg(j=h)\to
      \exists v\leq s\,
       (\mathsf E^{(2)}(s,\bar9,t,j,v)
          \land\mathsf E^{(2)}(s,\bar9,St,j,v))\bigr]\bigr].
\end{aligned}
\]
In particular, the selected row satisfies \(i<m\); an arbitrary
five-field value in the padding of family 1 cannot serve as a transition.
The left-move equation is \(h=Sh'\), so it cannot move left from shifted
column \(0\).  Set
\[
\begin{aligned}
 \mathsf H_3(s)\defeq{}&
  \exists k\leq s\,\exists n\leq s\,\exists b\leq s\,
  \exists L\leq s\,\exists T\leq s\,\exists H\leq s\,
  \exists W\leq s\,\exists m\leq s\,\exists\mu\leq s\,\bigl[\\
 &\quad\mathsf E(s,0,0,k)\land\mathsf E(s,0,\bar2,n)
       \land\mathsf E(s,0,\bar3,b)\land\mathsf E(s,0,\bar5,L)\\
 &\quad\land\mathsf E(s,0,\bar6,T)\land\mathsf E(s,0,\bar7,H)
       \land\mathsf E(s,0,\bar8,W)\\
 &\quad\land\mathsf{Le}(\bar1,L)\land H=T
       \land W=(L+\bar2\times T)+\bar1
       \land k=\bar3+m\land W=S\mu\\
 &\quad\land\bigl[\forall t\leq T\,\mathsf{Cfg}(s,t,n,b,W,\mu)\bigr]\\
 &\quad\land\mathsf{Init}(s,L,H,W)
       \land\forall t<T\,\mathsf{Step}(s,t,m,W)\bigr].
\end{aligned}
\]

For the output bits, their values and their occurrence on the final tape,
let
\[
\begin{aligned}
 \mathsf{OutCells}(s,T,H,\ell,y)\defeq{}&
  \forall i<\ell\,\exists d\leq s\,\exists v\leq s\,\bigl[
     \mathsf E(s,\overline{12},i,d)
      \land\mathsf E(s,\overline{13},i,v)\\
 &\hspace{24mm}\land\mathsf{Bit}(d)\land\mathsf{Le}(v,y)
      \land\mathsf E^{(2)}(s,\bar9,T,H+i,Sd)\bigr],\\
 \mathsf{OutNext}(s,\lambda)\defeq{}&
  \forall i<\lambda\,\exists u\leq s\,\exists v\leq s\,
    \exists d\leq s\,\bigl[
      \mathsf E(s,\overline{13},i,u)
        \land\mathsf E(s,\overline{13},Si,v)\\
 &\hspace{24mm}\land\mathsf E(s,\overline{12},Si,d)
        \land v=\bar2\times u+d\bigr].
\end{aligned}
\]
The final check is
\[
\begin{aligned}
 \mathsf H_4(s)\defeq{}&
  \exists T\leq s\,\exists H\leq s\,\exists W\leq s\,
  \exists\ell\leq s\,\exists y\leq s\,
  \exists\lambda\leq s\,\exists d\leq s\,\bigl[\\
 &\quad\mathsf E(s,0,\bar6,T)\land\mathsf E(s,0,\bar7,H)
       \land\mathsf E(s,0,\bar8,W)
       \land\mathsf E(s,0,\bar9,\ell)
       \land\mathsf E(s,0,\overline{10},y)\\
 &\quad\land\ell=S\lambda\land\mathsf{Lt}(H+\ell,W)
       \land\mathsf E(s,\bar7,T,\bar1)
       \land\mathsf E^{(2)}(s,\bar9,T,H+\ell,0)\\
 &\quad\land\mathsf E(s,\overline{12},0,d)
       \land\mathsf E(s,\overline{13},0,d)
       \land\mathsf E(s,\overline{13},\lambda,y)\\
 &\quad\land\bigl((y=0\land\ell=\bar1\land d=0)
                    \lor(\neg(y=0)\land d=\bar1)\bigr)\\
 &\quad\land\mathsf{OutCells}(s,T,H,\ell,y)
       \land\mathsf{OutNext}(s,\lambda)\bigr].
\end{aligned}
\]
Only the output prefix and its first following blank are constrained
by this final check.  The final head position, negative physical cells,
and cells beyond that blank impose no further output condition.
Define the promised computation formula by literal substitution of
the displayed abbreviations:
\[
 \mathsf{Comp}(e,x,s)\defeq
   \mathsf H_1(x,s)\land\mathsf H_2(e,s)
                   \land\mathsf H_3(s)\land\mathsf H_4(s).
\]
All occurrences of brackets denoting numerical entries in the proofs
below remain metatheoretic.  Every lookup inside the formulas above
is an \(\mathsf E\), \(\mathsf E^{(2)}\), or previously defined
\(\mathsf R_5\) abbreviation, with its displayed relational expansion.
\end{definition}

\begin{lemma}[Correctness of the configuration and output checks]
\label{lem:p04-configuration-output-formulas}
The formulas \(\mathsf H_3,\mathsf H_4\), and hence
\(\mathsf{Comp}\), are bounded arithmetic formulas.
Fix \(e,x,s\in\N\) satisfying \(\mathsf H_1(x,s)\land\mathsf H_2(e,s)\)
in \(\mathcal N\).  Then \(\mathsf H_3(s)\) holds if and only if the
header and tagged configuration families satisfy clause (h3), with
the fixed window of Definition~\ref{def:history-coding}, and describe
an initial segment of length \(T\) transitions of the raw decoded
machine on \(\operatorname{bin}(x)\).
If these conditions hold, \(\mathsf H_4(s)\) holds if and only if
the final configuration satisfies clause (h4) and the stored output
length, digits, prefixes and value are its canonical numerical output.
\textit{Status: complete proof.}
\end{lemma}

\begin{proof}
Boundedness follows by expansion.  Every quantifier in the new
definitions is marked bounded, every bound is a term of \(\Lang_A\),
and every comparison within a matrix is expressed by
\(\mathsf{Lt}\), \(\mathsf{Le}\), or an atomic equality and Boolean
connectives.  In particular, \(St\), \(Sj\), \(H+i\), and \(Sd\)
are terms, and both levels of each \(\mathsf E^{(2)}\) lookup have
the previously proved arithmetic bounds.  Fresh renaming before
substitution prevents capture.  No subtraction, division, power,
tuple constructor or tape-access function occurs in the expanded
formula.

By the H1 and H2 correctness lemmas, \(e\) decodes to the specified
valid machine, the input data are its canonical binary word, and
the header values \(k,n,b,L\) are fixed uniquely.  All further
header or configuration witnesses are likewise unique extracted
values.  If H3 holds, its equations give
\[
 H=T,\qquad W=L+2T+1,\qquad m=k-3=(n-2)b,\qquad \mu=W-1,
 \qquad L\geq1.
\]
Consequently every input column \(H+i\), \(i<L\), is below \(W\).

Write \(q_t=s[7,t]\), \(h_t=s[8,t]\), and
\(a_{t,j}=s[9,t,j]\) in the metatheory.  The formula Cfg gives
\(q_t<n\), \(h_t<W\), and \(a_{t,j}<b\) at every prescribed time
and column.  Its remaining conjuncts state the first-prefix value,
each base-\(b\) prefix recurrence, and equality of the last prefix
with \(s[11,t]\).  Induction over the columns therefore identifies
all stored tape prefixes and the complete fixed-width tape code.
These prefixes are nondecreasing because \(b\geq3\) and their
digits are nonnegative; each is at most the final tape code.
Leading zero tape digits are allowed, as the width is specified
separately.

Define a full tape in the metatheory by
\[
 \tau_t(z)=
 \begin{cases}
    a_{t,z+H},&-H\leq z<W-H,\\
    0,&\text{otherwise},
 \end{cases}
 \qquad z\in\Z,
\]
where \(0\) is the normalized blank symbol, and define its physical
head to be \(h_t-H\).  Its nonblank support is finite.
Init gives \(q_0=0\) and physical head \(0\).  It sets the columns
\(H+i\) to \(Sd_i\), the correct symbols for the input bits.
Every other column is blank: for \(j<W\), failure of
\(j<H\) and \(H+L\leq j\) is exactly \(j=H+i\) for a unique
\(i<L\).  Extending by blank outside the window therefore gives
precisely the original initial tape.

For \(t<T\), Step selects \(i<m\).  The certified-table lemma
therefore identifies its R5 tuple with an actual row of the decoded
transition table; padding rows are unavailable.  Its read symbol
is \(a_{t,h_t}\), its source state is \(q_t\), and its next state
is \(q_{t+1}\).  The explicit inequalities in Cfg put both heads
inside the window.  The two head equations give exactly a left
or right physical move.  The scanned column receives the row's
written symbol; the universal conjunct preserves every other
column.  Outside the window both full tapes are blank and the
scanned cell is inside, so those cells are unchanged as well.
Thus the full configurations have precisely the machine's
one-step relation.  Neither halting state occurs before time \(T\).
Induction on \(t\) proves that this is the actual initial run
segment, using determinism of the decoded table.

Conversely, suppose the tagged configurations and header have the
data in clause (h3), including the stated fixed-width certificates,
and describe that run segment.  Choose the actual stored header,
state, head, symbol and prefix values as witnesses.  They are at
most \(s\).  The auxiliary witnesses satisfy
\[
 m=k-3\leq k\leq s,\qquad \mu=W-1<W\leq s.
\]
Cfg and Init now follow from the stated ranges, recurrences and
initial tape.  At each transition the actual table row has a
unique index \(i<m\).  Its code is a stored description entry,
so the row fields and all auxiliary pairing tails are at most
that code, and hence at most \(s\).  These supply every Step
witness.  Its head equation, written-cell equation and unchanged-cell
conjuncts hold by the actual transition.  This proves H3.

Assume now these equivalent H3 conditions.  If H4 holds, its unique
header values are the same \(T,H,W\).  Put \(o_i=s[12,i]\) and
\(w_i=s[13,i]\).  The output clauses assert
\[
 \ell\geq1,\quad H+\ell<W,\quad q_T=1,\quad
 a_{T,H+\ell}=0,\quad a_{T,H+i}=So_i\quad(i<\ell).
\]
Thus all cells from physical \(0\) through \(\ell-1\) are
nonblank binary symbols and physical cell \(\ell\) is blank.
It is the first nonnegative blank.  The remaining clauses give
\(w_0=o_0\), \(w_{i+1}=2w_i+o_{i+1}\), \(w_{\ell-1}=y\),
with bits \(o_i\) and \(w_i\leq y\).  The same quotient/remainder
argument used for H1, with the explicit zero and first-bit clauses,
identifies this word as exactly \(\operatorname{bin}(y)\).
This proves (h4) and the output-certificate assertion.

Conversely suppose the final configuration and the stored output
data have those properties.  Every stored header, bit and prefix
is at most \(s\); the witness \(\lambda=\ell-1\) is at most
\(\ell\leq s\).  All OutCells and OutNext conjuncts hold from
the binary word and its prefixes, whose values are nondecreasing
and at most \(y\).  The first following blank and accepting state
give the other conjuncts of H4.  This proves the reverse direction.
The argument imposes no condition on cells outside the output prefix
except its first following blank, or on the final head.
\end{proof}

\begin{theorem}[Bounded verification of computations]
\label{thm:bounded-verification}
There is a bounded \(\Lang_A\)-formula \(\mathsf{Comp}(e,x,s)\) such that:
\begin{enumerate}[label=(\alph*)]
\item for all \(e,x,s\in\N\),
      \(\mathcal N\entails\mathsf{Comp}(\bar e,\bar x,\bar s)\) iff
      \(s\) is a history code for \((e,x)\), so that \(s\) codes a valid
      accepting computation of program \(e\) on input \(x\) in the fixed
      machine numbering whose final tape has a valid numerical output;
      and
\item \(\varphi_e(x)\) is defined iff
      \(\mathcal N\entails\exists s\,\mathsf{Comp}(\bar e,\bar x,s)\).
\end{enumerate}
The symbol \(\mathsf{Comp}\) stands
for this arithmetic formula, not for the metalinguistic relation it defines;
the distinction is used in the theorem and exercises immediately below.
\textit{Status: complete proof.}
\end{theorem}

\begin{proof}
Use exactly the formula of Definition~\ref{def:p04-computation-formula}.
Its boundedness was proved in the preceding lemma.  For (a), the H1
and H2 lemmas identify the input and program certificates with clauses
(h1) and (h2).  Under those conditions,
Lemma~\ref{lem:p04-configuration-output-formulas} identifies H3 and
H4 with (h3) and (h4), including the actual decoded transition
function and numerical-output convention.  All uses of a header
position agree by uniqueness of tagged extraction.  Thus the four
conjuncts hold exactly for the history codes in the statement.
If \(e\) is invalid, H2 has no witness, which also agrees with
the definition of a history code.

One direction of (b) now follows immediately from (a): a satisfying
code supplies a finite accepting computation of the raw decoded
machine with numerical output, so \(\varphi_e(x)\) is defined by
the fixed numbering.  For the converse, suppose \(\varphi_e(x)\)
is defined, with output \(y\).  Then \(e\) is valid and its raw
machine has a finite accepting computation, say of \(T\) transitions,
on the canonical word for \(x\).  Let \(L\geq1\) be the input length
and choose the numerical header values
\[
 H=T,\qquad W=L+2T+1.
\]
Write \(p_t\) for the physical head position at time \(t\).
Induction on \(t\), starting with \(p_0=0\) and using the two
possible moves, gives \(-t\leq p_t\leq t\).
Only the initial input cells \(0,\ldots,L-1\) and cells scanned
before time \(T\) can be nonblank at a time at most \(T\).
Their positions all lie in
\[
 [-T,\max(T,L-1)].
\]
The chosen columns \(0,\ldots,W-1\) represent the physical
interval \([-T,L+T]\), which contains this support and every
head position.  Moreover the rightmost physical cell \(L+T\)
is blank throughout: it is strictly larger than both \(T\)
and \(L-1\).  Set \(h_t=p_t+H\).  Then
\(0\leq h_t\leq2T<W\).  Both heads of each transition are
inside the window; a left move before time \(T\) cannot leave
it, since \(h_t\geq T-t>0\) for \(t<T\).

Store the actual states, shifted heads and tape symbols in the
prescribed families.  For every row of tape symbols, form its
successive base-\(b\) values, and store the last as its tag-11
complete tape code.  These are finite natural numbers, even
though their magnitudes need not have a polynomial bound in the
run length.  The last blank observation also bounds the first
nonnegative blank \(\ell\) of the final tape by \(L+T\).
Since the accepted output is canonical, \(\ell\geq1\), and
\[
 H+\ell\leq L+2T=W-1<W.
\]
Store its bits, numerical prefixes, length and value \(y\).
The binary prefix values are at most \(y\), by the same
monotonicity argument as for H1.

Take the finite input and table assignment from
Lemma~\ref{lem:p04-input-table-packing} and add these header,
configuration and output entries.  Shared header values agree,
and the family tags have the fixed positional arities of the
history definition, so the assignment is consistent.  Its
finitely many indices are computed from the actual run before
choosing \(s\).  The finite-packing lemma now supplies one \(s\)
realizing them all and bounds every stored value by \(s\).
All remaining existential witnesses meet their bounds:
\(m=k-3\leq k\leq s\), \(\mu=W-1<W\leq s\), and
the last output position is below \(\ell\leq s\);
the input last position was treated in the H1 proof.
Each transition uses an actual row \(i<m\), whose fields and
pairing tails are bounded by its stored row code.  The E and E2
formulas supply the bounds for their paired indices independently
of whether those indices are at most \(s\).
The two input/table lemmas and the configuration/output lemma
therefore give all four conjuncts of \(\mathsf{Comp}(e,x,s)\).
This proves the other direction of (b).

The argument includes the empty step range \(T=0\) for H3:
it describes just the initial configuration.  In the normalized
numbering the start state is \(0\), distinct from accepting state
\(1\), so H4 cannot hold for that zero-step segment.  Rejecting
halts and accepting halts with noncanonical output likewise fail
the specified final conjuncts.
\end{proof}

The formulas and proofs above concern the handbook's exact numerical
serialization and tagged history layout.  For other standard coding
constructions see \cite{Shoenfield1967,BoolosBurgessJeffrey2007}.
The result proved here is a truth equivalence in \(\mathcal N\);
numeralwise representation in \(\mathsf Q\) and the internal
proof transformations used for provability predicates require their
separate arguments in the incompleteness part.

\begin{example}[A one-step numerical identity computation]
Take \(n=b=3\), input-alphabet bitmap \(A=6\), and the three rows
\[
 \delta(0,a)=(1,a,R),\qquad
 \rho_a=\operatorname{code}_5(0,a,1,a,1)\quad(a=0,1,2).
\]
Put \(c_0=\operatorname{code}_6(3,3,6,\rho_0,\rho_1,\rho_2)\)
and \(e=\operatorname{pair}_{\N}(6,c_0)\).  This is a valid
description with \(m=3\) and \(k=6\): every nonhalting key occurs,
each transition preserves the scanned symbol and accepts after
one right move, and both binary symbols belong to the input alphabet.
For \(x=2\), the input bits are \((1,0)\), with prefixes \((1,2)\).
Use header
\[
 (k,c_0,n,b,A,L,T,H,W,\ell,y)
       =(6,c_0,3,3,6,2,1,1,5,2,2).
\]
The bitmap bits are \((0,1,1)\), with tails \((6,3,1,0)\).
At both times \(0\) and \(1\), the five tape symbols are
\[
 (0,2,1,0,0),
\]
whose base-three prefixes are \((0,2,7,21,63)\); both complete
tape codes are \(63\).  The states are \((0,1)\) and the shifted
heads are \((1,2)\).  The step reads symbol \(2\) at column \(1\)
and uses row \(i=2<m\), leaving that symbol unchanged.
The output bits and prefixes are again \((1,0)\) and \((1,2)\),
with the first following blank at column \(H+\ell=3<5\).
Together with the unique program-decoding tails, these data form
a compatible finite assignment.  Beta packing supplies an \(s\)
for which \(\mathsf{Comp}(e,2,s)\) holds.  Printing the large
decimal value of that packing code is unnecessary for the exact check.

For a failure case, the word \((0,1)\), with prefixes \((0,1)\),
has value \(1\) and satisfies the digit recurrence and prefix bounds.
It nevertheless fails H1 for \(x=1\), because its first digit is
zero.  Thus the recurrence alone would accept a noncanonical input.
Similarly, the bitmap \((1,1,1)\), with value \(7\), fails H2:
its bit at the normalized blank symbol must be zero.
\end{example}

\begin{theorem}
\label{thm:p04-ce-sigma-one}
A set is computably enumerable iff it belongs to
\(\boldsymbol{\Sigma}^0_1\).  A set is decidable iff it belongs to
\(\boldsymbol{\Delta}^0_1\).
\textit{Status: complete proof.}
\end{theorem}

\begin{proof}
First suppose \(A\) is computably enumerable.  Its defining
domain characterization in Section~\ref{sec:turing-machines-03}
gives a partial computable function with domain \(A\).
Theorem~\ref{thm:enumeration} supplies an index \(e\) for that function
in this exact numbering.  Bounded verification now gives
\[
 x\in A\quad\Longleftrightarrow\quad
 \mathcal N\entails\exists s\,\mathsf{Comp}(\bar e,\bar x,s).
\]
Substituting the fixed numeral \(\bar e\) does not change
boundedness of the matrix.  Its fixed negation-normal-form
conversion remains bounded, so the displayed formula is
\(\Sigma^0_1\) in the declared syntactic convention.  This proves
\(A\in\boldsymbol{\Sigma}^0_1\).

For the converse, first verify that the truth of each fixed bounded
arithmetic formula on a supplied finite tuple of numbers is decidable.
Evaluate terms recursively: \(0\), variables and successor are
immediate, addition is repeated increment, and multiplication is
repeated addition.  Return false for \(\bottom\).  Compare term values
for an atomic equality, and evaluate Boolean connectives by their
truth tables.  For a marked bounded quantifier, evaluate its bound
under the current assignment and then evaluate its matrix on each
number in the finite specified range.  The quantified variable does
not occur in that bound, so the range is fixed for this loop.
An existential returns true when a matrix value is true, and
otherwise false after exhaustion; a universal returns false on
the first false value, and otherwise true.  This also gives the
correct false/true answers for an empty existential/universal range.
Structural induction on the formula proves termination and correctness
of this evaluator, since every recursive call concerns a proper
subformula and each loop has finitely many iterations.

This is a machine procedure on finite numerical words: retain the
assignment and loop counters in delimited blocks, evaluate a fixed
term or subformula in scratch blocks, and use scans, copying and
binary increment for the indicated arithmetic loops and comparisons.
The formula is fixed, so its finite syntax and the return points
for its nested loops are part of finite control; the counters and
intermediate values occupy finite, extendible tape blocks.
Each scan or increment terminates on its delimited finite block,
and the structural-induction argument proves that the complete
procedure halts.  To return its truth value, put the corresponding
single binary digit at physical cell \(0\), put a blank at cell \(1\),
and enter the accepting state; later scratch cells are immaterial to
the output convention.  This computes the required characteristic function.
The construction
uses neither arithmetical representability nor the theorem being proved.

Now let a \(\Sigma^0_1\) formula define \(A\).
Theorem~\ref{thm:prefix-normal-form} gives an \(\mathcal N\)-equivalent
formula
\[
 \exists z_1\cdots\exists z_r\,\theta(x,z_1,\ldots,z_r),
\]
with bounded \(\theta\).  A vacuous existential can be added if
needed, since \(\N\) is nonempty.  On input \(x\), for stages
\(v=0,1,2,\ldots\), enumerate all \(r\)-tuples with every coordinate
at most \(v\), using finite nested counters, and apply the total
bounded-formula evaluator to each tuple.  If any value is true,
write the canonical output \(0\) and accept.  Each stage terminates,
and every possible witness tuple appears at a finite stage.
Thus this machine computes the partial function equal to \(0\)
on \(A\) and undefined elsewhere.  If there is no witness it
continues through infinitely many stages, so it does not halt
with an unintended output.  The domain definition proves that
\(A\) is computably enumerable, including when \(A\) is empty.

Finally, negation interchanges the \(\Sigma^0_1\) and
\(\Pi^0_1\) definability classes: in prefix form,
the negation of \(\forall\bar z\,\theta\) is
\(\exists\bar z\,\neg\theta\), with bounded matrix, and conversely.
Consequently \(A\in\boldsymbol{\Delta}^0_1\) if and only if
both \(A\) and its complement are \(\Sigma^0_1\)-definable.
By the first assertion this holds exactly when both sets are
computably enumerable, which is the earlier characterization
of decidability in Section~\ref{sec:turing-machines-03}.
This proves the second assertion.
\end{proof}

\begin{definition}[Many-one hardness and completeness for a class]
Let \(\mathcal C\) be a class of subsets of \(\N\).  A set \(B\) is
\emph{\(\mathcal C\)-hard under many-one reductions} iff
\(A\leq_m B\) for every \(A\in\mathcal C\).  It is
\emph{\(\mathcal C\)-complete under many-one reductions} iff, in addition,
\(B\in\mathcal C\).  This term is required to distinguish an upper-bound
classification from the uniform lower bound used in the next theorem.  It is
a property of a set relative to \(\mathcal C\) and \(\leq_m\), and is not
semantic completeness, deductive completeness, or the status of a proof.
The theorem and the following exercise use exactly this distinction.
\end{definition}

\begin{theorem}[Hierarchy classifications of program domains and totality]
The set \(K\) is \(\boldsymbol{\Sigma}^0_1\)-complete under many-one
reductions, and its complement is \(\boldsymbol{\Pi}^0_1\)-complete under
many-one reductions.  The totality set
\[
  \mathsf{TOT}=\set{e:\forall x\,\exists k\,
  \text{program }e\text{ accepts on }x\text{ with a valid numerical output}
  \text{ within }k\text{ steps}}
\]
is \(\boldsymbol{\Pi}^0_2\)-complete under many-one reductions.
\textit{Status: cited theorem.}
\end{theorem}

See \cite{Shoenfield1967} for complete proofs.  The displayed prefix gives
the upper bound for \(\mathsf{TOT}\), read into the closure classes by
\cref{thm:prefix-normal-form}; the completeness conclusion additionally
requires a uniform many-one reduction.  The cited completeness proofs
classify relations by prefix-form formulas, and
\cref{thm:prefix-normal-form} licenses their transfer to the classes
fixed here.

\begin{theorem}[Strictness of the arithmetical hierarchy]
For every \(n\geq1\),
\[
  \boldsymbol{\Delta}^0_n\subsetneq\boldsymbol{\Sigma}^0_n,
  \qquad
  \boldsymbol{\Delta}^0_n\subsetneq\boldsymbol{\Pi}^0_n,
  \qquad
  \boldsymbol{\Sigma}^0_n\neq\boldsymbol{\Pi}^0_n,
\]
and
\(\boldsymbol{\Sigma}^0_n\cup\boldsymbol{\Pi}^0_n
\subsetneq\boldsymbol{\Delta}^0_{n+1}\), for unary relations.
\textit{Status: cited theorem.}
\end{theorem}

See \cite{Shoenfield1967} for a complete proof.  The cited proof works
with prefix-form classes; \cref{thm:prefix-normal-form} transfers the
statement to the closure classes fixed here.

\begin{warning}[Syntactic prefix versus semantic class]
The syntactic classification of a displayed formula supplies an upper bound
on the definability class of the relation it defines.  A completeness
reduction or a nondefinability argument is needed to show that the relation
does not lie in a lower boldface class.
\end{warning}
